\documentclass[conference]{IEEEtran}
\IEEEoverridecommandlockouts

\usepackage{cite}
\usepackage{amsmath,amssymb,amsfonts}
\usepackage{algorithmic}
\usepackage{graphicx}
\usepackage{textcomp}
\usepackage{xcolor}
\usepackage{booktabs}
\usepackage{multirow}
\usepackage{listings}
\usepackage{url}

% Placeholder box for figures (replace with \includegraphics when adding images)
\newcommand{\figplaceholder}[1]{\fbox{\parbox[c][3.2cm][c]{0.92\columnwidth}{\centering\footnotesize Figure placeholder:\\ \texttt{#1}}}}

\lstset{basicstyle=\ttfamily\scriptsize, breaklines=true, frame=single, columns=fullflexible,
        keywordstyle=\bfseries, language=Python, showstringspaces=false}

\begin{document}

\title{Transformer-Based Real-Time Adaptive Quantum Error Mitigation for Non-Stationary Noise}

\author{\IEEEauthorblockN{Madhu Hasini}
\IEEEauthorblockA{\textit{CB.AI.U4AID24060} \\
\textit{[Department, Institution]}\\
[City, Country] \\
[email]}
\and
\IEEEauthorblockN{Saanvika}
\IEEEauthorblockA{\textit{CB.AI.U4AID24047} \\
\textit{[Department, Institution]}\\
[City, Country] \\
[email]}
\and
\IEEEauthorblockN{Dharshini}
\IEEEauthorblockA{\textit{CB.AI.U4AID24013} \\
\textit{[Department, Institution]}\\
[City, Country] \\
[email]}
}

\maketitle

\begin{abstract}
Quantum error mitigation methods such as zero-noise extrapolation and Pauli twirling are usually run with settings fixed in advance, even though the noise of superconducting processors changes from day to day. This project asks whether forecasting that change with a Transformer leads to better mitigation decisions than three alternatives: a fixed configuration, adapting to the latest published calibration, and a contextual multi-armed bandit, which is the approach of our base paper. We built a leakage-safe dataset of 101,891 real calibration states covering 600 qubits on four IBM Heron processors, trained per-qubit and per-pair forecasters, and used the forecasts to choose among 14 mitigation actions for six benchmark circuits. Circuit execution was simulated with noise models built from the real calibrations, and the test period was evaluated once after all choices were frozen. The Transformer forecast calibrations 10 to 26 percent better than persistence but no better than a window mean, ridge regression or gradient-boosted trees. Choosing the mitigation from these forecasts did not beat choosing it from the latest calibration at horizons of one, four or seven days. Adapting to the calibration itself did help: error fell by 4.2 to 5.6 percent relative to the best fixed configuration and by 15 to 19 percent relative to the bandit at an equal shot budget. The largest possible benefit of a perfect forecast was 1.4 to 7.1 percent and was concentrated in abrupt device changes.
\end{abstract}

\begin{IEEEkeywords}
quantum error mitigation, zero-noise extrapolation, Pauli twirling, calibration drift, Transformer, contextual bandit, IBM Quantum
\end{IEEEkeywords}

%=====================================================================
\section{Introduction}
%=====================================================================
Present quantum processors are noisy, and full error correction is still out of reach for most applications \cite{preskill2018}. Quantum error mitigation (QEM) is the practical middle ground: instead of correcting errors during the computation, it post-processes the outputs of several related circuit runs to estimate what a noise-free device would have produced \cite{cai2023}. Zero-noise extrapolation (ZNE) is one of the most widely used of these methods. It runs a circuit at several amplified noise levels and extrapolates the measured expectation value back to zero noise \cite{temme2017,li2017,kandala2019,giurgica2020}. Pauli twirling, also known as randomized compiling, converts coherent gate errors into stochastic Pauli errors, which makes the noise easier to model and extrapolate \cite{wallman2016}.

ZNE has several settings that the user must pick: which noise-scale factors to run, which extrapolation function to fit, and whether to twirl. These settings are normally chosen once. The noise they are meant to handle, however, is not constant. Relaxation times, gate errors and readout errors of superconducting qubits drift from one calibration to the next and sometimes jump abruptly \cite{klimov2018,proctor2020}. A setting that works well on one day can be a poor choice a few days later.

Our base paper, by Rahman and Nguyen \cite{rahman2026}, addresses this with a contextual multi-armed bandit that picks the ZNE folding level online. Their method reacts to recent feedback and was tested against a prescribed, synthetic drift profile. In this project we ask a different question: can we look ahead? If a Transformer can forecast how a device's calibration will change, the mitigation could be configured for the noise that will be present when the circuit actually runs, rather than for the noise that was measured earlier.

We set out to test this carefully. We pre-registered our hypotheses and success criteria before running the main experiments, used real IBM calibration histories instead of synthetic drift, compared against strong simple baselines as well as the bandit, and evaluated the test period only once. The answer turned out to be more nuanced than we expected, and we report it as it is.

%=====================================================================
\section{Related Work}
%=====================================================================
\textit{Zero-noise extrapolation and twirling.} Temme et al. \cite{temme2017} and Li and Benjamin \cite{li2017} introduced the idea of extrapolating to zero noise. Kandala et al. \cite{kandala2019} demonstrated it on superconducting hardware by stretching pulses, and Giurgica-Tiron et al. \cite{giurgica2020} showed that unitary folding achieves the same effect digitally. Kim et al. \cite{kim2023} combined ZNE with Pauli twirling at a scale of over one hundred qubits. Wallman and Emerson \cite{wallman2016} showed that twirling tailors arbitrary noise into a Pauli channel. Cai et al. \cite{cai2023} give a broad review of these and other mitigation techniques.

\textit{Adaptive and learning-based mitigation.} Rahman and Nguyen \cite{rahman2026} select the ZNE folding set with a cost-aware LinUCB contextual bandit \cite{li2010} under time-varying noise and report up to 40 percent fewer circuit executions than fixed folding. Hour et al. \cite{hour2026} adjust the folding to the device noise. Learning-based estimators such as Clifford data regression \cite{czarnik2021,perezguijarro2026} learn a correction from near-Clifford training circuits, and Liao et al. \cite{liao2024} apply machine learning directly to mitigation. Placidi et al. \cite{placidi2026} train attention-based sequence models that map noisy output distributions to ideal ones. Prodius et al. \cite{prodius2026} study how finite-shot uncertainty propagates through mitigation estimates, which is relevant to choosing between low-variance and low-bias settings.

\textit{Drift in quantum hardware.} Proctor et al. \cite{proctor2020} developed methods to detect and track drift in quantum processors, and Klimov et al. \cite{klimov2018} documented large fluctuations of energy-relaxation times caused by two-level-system defects. Neither line of work uses forecasts of drift to configure mitigation.

\textit{Sequence models.} The Transformer \cite{vaswani2017} relies on self-attention to model dependencies in sequences and is a natural candidate for forecasting calibration histories.

%=====================================================================
\section{Literature Gaps Identified}
%=====================================================================
\begin{itemize}
\item \textbf{No forecasting.} Adaptive ZNE in the base paper \cite{rahman2026} is reactive. It does not try to predict the noise at the time of execution.
\item \textbf{Synthetic drift.} The base paper evaluates under a prescribed drift profile, not under measured calibration histories of real processors.
\item \textbf{Narrow action space.} Only the folding set is adapted. The extrapolation function and Pauli twirling stay fixed.
\item \textbf{Hidden cost of feedback.} A bandit needs a reward. For classification the reward comes from labels; when estimating expectation values of arbitrary circuits, the ideal value is unknown, and probe circuits cost shots.
\item \textbf{Weak baselines.} Learned forecasters are rarely compared with simple ones such as the latest value or a moving average, and forecast accuracy is rarely linked to the quality of the decisions it drives.
\end{itemize}

%=====================================================================
\section{Why Quantum}
%=====================================================================
The problem only exists because of quantum hardware. Mitigation settings matter because noisy gates on real superconducting qubits bias expectation values, and the best setting depends on physical quantities ($T_1$, $T_2$, single- and two-qubit gate errors) that change as the device's defect landscape and calibration change \cite{klimov2018}. Expectation values of entangled multi-qubit states cannot be read off a classical model of the circuit at scale, which is why mitigation and not classical simulation is used on hardware \cite{kim2023}. Our study uses real calibration records of IBM Heron processors and evaluates mitigation on circuits whose ideal values are known, so that the effect of every decision can be measured exactly.

%=====================================================================
\section{Our Contributions}
%=====================================================================
\begin{itemize}
\item A leakage-safe benchmark that links real IBM Heron calibration drift (101,891 calibration states, 600 qubits, four processors) to adaptive ZNE and Pauli-twirling decisions, with chronological, unseen-qubit and unseen-device evaluation.
\item A calibration-informed, model-based controller that jointly selects ZNE scale factors, extrapolator and twirling, and that reduces error by 4.2 to 5.6 percent against the best fixed configuration and by 15 to 19 percent against a faithful contextual-bandit baseline at an equal shot budget.
\item A pre-registered test of Transformer-based forecasting showing that better calibration forecasts did not produce better mitigation decisions, together with a value-of-information analysis that bounds the possible benefit at 1.4 to 7.1 percent and locates it in abrupt device changes.
\end{itemize}

%=====================================================================
\section{Methodology}
%=====================================================================
\subsection{Dataset}
We used the public \textit{qiskit-calibration-drift} dataset \cite{driftdata}, which records the published calibration properties of IBM Heron processors every 30 minutes together with weather and space-weather indicators. We pinned a single revision and verified its checksum. The original project used only the first 1,000 rows. We used the full file and filtered it carefully (Table~\ref{tab:data}).

\begin{table}[htbp]
\caption{Calibration data used in this study (real IBM records)}
\begin{center}
\begin{tabular}{lr}
\toprule
\textbf{Item} & \textbf{Value} \\
\midrule
Rows in the dataset file & 16,140,252 \\
Genuine measurements (value changed) & 2,024,027 \\
Calibration states used for modelling & 101,891 \\
Backends & fez, kingston, marrakesh, torino \\
Qubits & 600 \\
Period & Jan. to Oct. 2026 \\
Median spacing between calibrations & 24.8 h \\
Properties & $T_1$, $T_2$, sx error, readout, CZ error \\
\bottomrule
\end{tabular}
\label{tab:data}
\end{center}
\end{table}

\textit{Cleaning.} $T_1$ and $T_2$ were stored in seconds before a schema change in May 2026 and in microseconds after it; we converted everything to microseconds. This mix of units had also affected the first version of our pipeline. An event counts as a genuine measurement only if its value differs from the previous event of the same property and qubit; about 89 percent of the single-qubit gate-error events were re-timestamps of unchanged values. Failure ceilings (values at or above 0.999) were removed. Every other property is attached to a calibration state by an as-of join on its publication time, so a state never contains information that was not public at the decision time. A gap of more than seven days starts a new segment, and windows never cross a segment.

\textit{Splits.} The split is chronological: training before 1 August 2026, validation in August, and test from 1 September to 12 October 2026. ibm\_fez and ibm\_kingston are the development backends. Twenty percent of their qubits are held out entirely, ibm\_marrakesh is held out as an unseen device, and the retired ibm\_torino is used only descriptively. Scalers and hyperparameters were fitted on training and validation data only.

\subsection{Forecasting models}
Each sample is a window of the last 20 calibration states of a qubit (19 features) and the targets are the changes in log $T_1$, log $T_2$, log single-qubit error, log readout error and log mean CZ error at 1, 2, 4 and 7 calibrations ahead. We compared persistence (the latest value), the mean of the last 5 and the last 20 values, ridge regression and histogram gradient-boosted trees on lag features, an LSTM, and a Transformer encoder with learned positional embeddings and a multi-horizon output head. Hyperparameters were chosen on validation. The selected Transformer has one encoder layer, model width 32, four heads and a feed-forward width of 64, and we report a five-seed ensemble. Because the mitigation decision depends mostly on two-qubit gate error, we also trained the same model families on pair-level CZ error, at horizons of 1, 4 and 7 days.

\subsection{Backend model and circuits}
IBM recalibrates $T_1$ on a whole chip within seconds once per day, so one decision is made per backend-day. Decisions are made for chains of four physically coupled qubits taken from the real coupling map (35 development chains and 16 on the unseen device). For each chain and day we build a Qiskit Aer \cite{qiskit2024} noise model from the published calibration. Single-qubit gates carry thermal relaxation over the real gate duration (24, 32 or 36 ns), and CZ gates carry relaxation over 68 ns. A depolarizing channel is added so that the average gate infidelity matches the reported error (within 0.07 percent in our checks). In a coherent regime we add a fixed controlled-phase error of 0.03 rad to every CZ. We used six circuits on the native gate set \{CZ, sx, rz, x\}: a hardware-efficient ansatz, a Trotterised transverse-field Ising model and a QAOA circuit, each at two depths (12 to 36 CZ gates). The observables are single-qubit and nearest-neighbour Z and X Pauli strings whose noiseless value exceeds 0.2 in magnitude.

\subsection{Mitigation actions and policies}
All policies choose from the same 14 actions: no mitigation, twirling only, and six ZNE settings with and without twirling (scale sets \{1,3\}, \{1,2,3\}, \{1,3,5\} and \{1,3,5,7\}; linear, Richardson, exponential and quadratic fits). Noise is scaled by local CZ folding, and twirling uses eight random Pauli frames. Every policy spends the same budget of 8,192 shots per observable.

The \textit{model-based controller} builds a noise model from a state estimate, simulates every action, and picks the one with the lowest expected error including shot noise. The state estimate is what distinguishes the policies: the latest published calibration (persistence), the window-mean, ridge or Transformer forecast, or the true execution state (an unrealisable oracle). The \textit{best fixed} policy uses the single action with the lowest validation error. The \textit{contextual bandit} follows the LinUCB formulation of the base paper \cite{rahman2026}: the reward is the error on a Clifford version of the circuit, whose ideal value is known, run with one eighth of the shot budget, and feedback arrives only after execution. Its hyperparameters were tuned on validation, and we added reward centering and a play-each-arm-once start, without which even a bandit rewarded with the true error performed worse than a fixed action.

\subsection{Metrics and evaluation protocol}
Forecasts are scored by the mean absolute error of the log value, averaged per qubit so that each qubit counts once. A mitigation decision is scored by the mean over observables of the expected absolute error $\mathbb{E}\,|\hat{E}-E_{\mathrm{ideal}}|$, with shot noise averaged over 200 Monte Carlo trials using common random numbers for all policies. Comparisons are paired differences with 95 percent confidence intervals from a cluster bootstrap (10,000 resamples of qubits or chains within each backend), and Holm correction is applied within each family of tests. Our hypotheses were fixed before the expanded experiments: H1, the Transformer forecasts better than every baseline; H2 (primary), Transformer-driven adaptation beats persistence-driven adaptation one day ahead; H3, it beats the best fixed action and is not worse under stationary noise; H4, it beats the bandit. The test period was evaluated once, after the hashes of all inputs were frozen.

%=====================================================================
\section{Results}
%=====================================================================
\subsection{Forecasting calibration drift}
Table~\ref{tab:fc} shows the test error of the per-qubit forecasters. Every learned model is clearly better than persistence, by 9.7 to 25.7 percent depending on the property and horizon. The Transformer is not better than the simple models, though. For $T_1$ and $T_2$ a plain mean of the last 20 values does as well, and for readout error ridge regression is significantly better. None of the 15 pre-registered H1 cells passed, so \textbf{H1 is not supported}. The reason is visible in the data: about half of the next-day $T_1$ values differ by more than 25 percent from the current value, so most of the change behaves like noise around a slowly moving level, and averaging captures most of what can be predicted.

\begin{table}[htbp]
\caption{Per-qubit forecast error on the test set (MAE of log value; 1 / 7 calibrations ahead)}
\begin{center}
\setlength{\tabcolsep}{3pt}
\begin{tabular}{lcccc}
\toprule
\textbf{Property} & \textbf{Persist.} & \textbf{Win. mean} & \textbf{Ridge} & \textbf{Transformer} \\
\midrule
$T_1$ & 0.307 / 0.311 & 0.229 / 0.234 & 0.233 / 0.235 & 0.229 / 0.232 \\
$T_2$ & 0.236 / 0.241 & 0.178 / 0.182 & 0.181 / 0.184 & 0.180 / 0.183 \\
sx error & 0.241 / 0.326 & 0.243 / 0.252 & 0.222 / 0.245 & 0.216 / 0.242 \\
Readout & 0.371 / 0.389 & 0.292 / 0.309 & 0.287 / 0.303 & 0.294 / 0.314 \\
CZ (mean) & 0.192 / 0.246 & 0.191 / 0.207 & 0.183 / 0.201 & 0.173 / 0.200 \\
\bottomrule
\end{tabular}
\label{tab:fc}
\end{center}
\end{table}

% ---- Figure: forecast error vs horizon (add your image) ----
\begin{figure}[htbp]
\centering
\figplaceholder{results/stage3/figures/test/mae\_vs\_horizon\_test\_primary.pdf}
% \includegraphics[width=\columnwidth]{mae_vs_horizon_test_primary.pdf}
\caption{Forecast error versus horizon for each calibration property (test, development qubits, 95\% confidence intervals).}
\label{fig:fc}
\end{figure}

The pair-level CZ forecasters behaved similarly (Table~\ref{tab:pair}). They improved on persistence by 11 to 12 percent one day ahead and by 18 to 20 percent four and seven days ahead, with the Transformer level with ridge and gradient-boosted trees.

\begin{table}[htbp]
\caption{Pair-level CZ error forecasting on the test set (reduction in MAE relative to persistence, development pairs)}
\begin{center}
\begin{tabular}{lccc}
\toprule
\textbf{Model} & \textbf{1 day} & \textbf{4 days} & \textbf{7 days} \\
\midrule
Window mean & 3.5\% & 14.9\% & 14.2\% \\
Ridge & 11.0\% & 20.0\% & 19.2\% \\
Gradient-boosted trees & 12.2\% & 19.2\% & 18.0\% \\
Transformer & 11.5\% & 18.7\% & 18.1\% \\
\bottomrule
\end{tabular}
\label{tab:pair}
\end{center}
\end{table}

\subsection{Adaptive mitigation}
Table~\ref{tab:hyp} summarises the pre-registered tests on the test period, and Table~\ref{tab:err} gives the mean errors. The primary hypothesis H2 failed. Choosing the mitigation from the Transformer forecast changed the error by +0.02 percent relative to choosing it from the latest calibration, with a confidence interval that includes zero. The same holds four and seven days ahead and on the unseen device, and the window-mean and ridge forecasts made no difference either. Adapting to the calibration, on the other hand, clearly helped. It reduced the error by 5.6 percent one day ahead and 4.2 percent four days ahead relative to the best fixed action (H3), and it was not worse under stationary noise. Against the contextual bandit, the reduction was 15 to 19 percent at every horizon (H4).

\begin{table}[htbp]
\caption{Pre-registered hypothesis tests (test set, development chains, real drift)}
\begin{center}
\setlength{\tabcolsep}{3pt}
\begin{tabular}{p{2.2cm}p{3.6cm}p{1.6cm}}
\toprule
\textbf{Hypothesis} & \textbf{Result (95\% CI)} & \textbf{Verdict} \\
\midrule
H1 forecasting & Transformer never beats all baselines (0/15 cells) & Not supported \\
H2 vs persistence (1 day, primary) & $+0.02\%$ [$-0.29$, $+0.54$] & Not supported \\
H2 at 4 / 7 days & $-0.41\%$ / $-0.05\%$, n.s. & Not supported \\
H3 vs best fixed & $-5.6\%$ (1 d), $-4.2\%$ (4 d), $-0.3\%$ (7 d, n.s.) & Supported (1, 4 d) \\
H3 stationary & $-4.9\%$ [$-7.4$, $-1.0$], margin $+2\%$ & Supported \\
H4 vs bandit & $-15.0\%$ / $-18.7\%$ / $-14.5\%$ & Supported \\
\bottomrule
\end{tabular}
\label{tab:hyp}
\end{center}
\end{table}

\begin{table*}[htbp]
\caption{Mean absolute error per observable on the test set (development chains). R1: incoherent drift; R2: with coherent CZ phase error. The oracle uses the true execution state and is not realisable.}
\begin{center}
\begin{tabular}{llcccccccc}
\toprule
\textbf{Horizon} & \textbf{Regime} & \textbf{None} & \textbf{Best fixed} & \textbf{Bandit} & \textbf{Persistence} & \textbf{Window mean} & \textbf{Ridge} & \textbf{Transformer} & \textbf{Oracle} \\
\midrule
1 day & R1 & 0.0458 & 0.0247 & 0.0277 & 0.0239 & 0.0241 & 0.0240 & 0.0239 & 0.0235 \\
1 day & R2 & 0.0467 & 0.0258 & 0.0283 & 0.0236 & 0.0240 & 0.0237 & 0.0237 & 0.0233 \\
4 days & R1 & 0.0489 & 0.0263 & 0.0309 & 0.0259 & 0.0260 & 0.0259 & 0.0258 & 0.0251 \\
4 days & R2 & 0.0497 & 0.0275 & 0.0324 & 0.0258 & 0.0259 & 0.0258 & 0.0257 & 0.0248 \\
7 days & R1 & 0.0507 & 0.0267 & 0.0316 & 0.0272 & 0.0270 & 0.0271 & 0.0270 & 0.0258 \\
7 days & R2 & 0.0514 & 0.0281 & 0.0323 & 0.0274 & 0.0271 & 0.0273 & 0.0275 & 0.0255 \\
\bottomrule
\end{tabular}
\label{tab:err}
\end{center}
\end{table*}

% ---- Figure: hypothesis forest plot (add your image) ----
\begin{figure}[htbp]
\centering
\figplaceholder{results/stage5/figures/hypotheses\_test.pdf}
% \includegraphics[width=\columnwidth]{hypotheses_test.pdf}
\caption{Paired error differences of Transformer-driven adaptation against each alternative (test set, 95\% cluster-bootstrap confidence intervals).}
\label{fig:hyp}
\end{figure}

\subsection{Why forecasting did not help}
The oracle row in Table~\ref{tab:err} shows how much a perfect forecast could add over the latest calibration: 1.4 to 1.7 percent one day ahead, 3.0 to 3.7 percent four days ahead and 5.2 to 7.1 percent seven days ahead. When we sorted the test decisions by how much the chain's CZ error actually changed, almost all of this benefit came from the quarter with the largest changes (a median change of about 2.3 times). There the oracle beat persistence by 6.4, 9.3 and 14.1 percent at the three horizons, and adapting to the stale calibration was worse than not adapting at all. The Transformer forecasts were closer to the truth than persistence in 90 to 97 percent of these cases, but they recovered only about 1.2 to 1.4 percent at one and four days and nothing at seven days. Elsewhere they were slightly worse. The forecasts improve accuracy mostly by pulling values toward a recent average, and that does not change which mitigation action looks best. Forecast accuracy and decision value are therefore not the same thing.

% ---- Figure: error by horizon (add your image) ----
\begin{figure}[htbp]
\centering
\figplaceholder{results/stage5/figures/error\_by\_horizon\_test.pdf}
% \includegraphics[width=\columnwidth]{error_by_horizon_test.pdf}
\caption{Mean mitigated error by decision horizon for each policy (test set).}
\label{fig:hor}
\end{figure}

\subsection{Why the bandit underperformed}
The bandit was worse than the best fixed action, even after tuning. The reason is its feedback. On the test period the ranking of actions by the Clifford-probe reward agreed only weakly with the ranking by true error (median Spearman correlation between 0.03 and 0.56 depending on the circuit), and the probe picked the truly best action in only 6 to 14 percent of states. The probes also used 12.5 percent of the shots, and with feedback delayed by several days the bandit degraded further. Even with an unrealisable reward equal to the true error, it only matched a well-chosen fixed action per circuit. The base paper's classification task provides a reward through labels; estimating an expectation value of an arbitrary circuit does not.

% ---- Figure: probe reward quality (add your image) ----
\begin{figure}[htbp]
\centering
\figplaceholder{results/stage5/figures/probe\_reward\_quality.pdf}
% \includegraphics[width=\columnwidth]{probe_reward_quality.pdf}
\caption{Agreement between the bandit's Clifford-probe reward and the true error ranking of the 14 actions, per circuit.}
\label{fig:probe}
\end{figure}

\subsection{Robustness and an exploratory controller}
Two robustness checks, defined and frozen before they were run, gave the same picture (Table~\ref{tab:rob}): eight new random instances of the circuits, and chains of six qubits.

\begin{table}[htbp]
\caption{Robustness on the test set (Transformer-driven adaptation, 1 day ahead)}
\begin{center}
\setlength{\tabcolsep}{3pt}
\begin{tabular}{lccc}
\toprule
\textbf{Compared with} & \textbf{Main} & \textbf{New circuits} & \textbf{6 qubits} \\
\midrule
Persistence & $+0.02\%$ & $-0.42\%$ & $-0.35\%$ \\
Best fixed & $-5.6\%^{*}$ & $-5.7\%^{*}$ & $-3.6\%$ \\
Bandit & $-15.0\%^{*}$ & $-17.8\%^{*}$ & -- \\
Oracle gap & 1.4--1.7\% & 2.5\% & 2.6\% \\
\bottomrule
\multicolumn{4}{l}{$^{*}$ 95\% confidence interval excludes zero.}
\end{tabular}
\label{tab:rob}
\end{center}
\end{table}

Finally, we tried a decision-focused Transformer that reads the chain's calibration history and directly corrects the predicted error of each action. On validation it was 0.28 percent worse than persistence one day ahead, and a control model without any history did slightly better than both, which points to a fixed correction of the simulator rather than learned drift. We stopped this experiment before the test stage and report it only as exploratory work.

%=====================================================================
\section{Conclusion and Outlook}
%=====================================================================
We tested whether Transformer forecasts of real calibration drift can improve adaptive zero-noise extrapolation and Pauli twirling. On real IBM Heron calibration histories and a strict, pre-registered protocol, they did not: neither as forecasts handed to a controller nor as a decision-focused controller did the Transformer beat simply adapting to the latest calibration. What did help was adapting at all. A calibration-informed controller lowered the error by 4 to 6 percent against the best fixed configuration and by 15 to 19 percent against the contextual-bandit approach of our base paper, without extra shots. Most of the remaining benefit that a forecast could deliver lies in abrupt device changes, and forecasters that are more accurate on average did not capture it.

Several limits apply. Circuit execution was simulated, readout and idle errors were not modelled, chains had four or six qubits, and the next published calibration stood in for the true state during execution. The controller assumes a reasonably accurate noise model, and its advantage shrank when a coherent error was left out of the model. Future work should train forecasters for decision quality rather than accuracy, use them only when a large change is expected, test on newer data that has not been seen, and confirm the results on hardware.

%=====================================================================
\begin{thebibliography}{00}
\bibitem{preskill2018} J. Preskill, ``Quantum computing in the NISQ era and beyond,'' \textit{Quantum}, vol. 2, p. 79, 2018.
\bibitem{cai2023} Z. Cai, R. Babbush, S. C. Benjamin, S. Endo, W. J. Huggins, Y. Li, J. R. McClean, and T. E. O'Brien, ``Quantum error mitigation,'' \textit{Rev. Mod. Phys.}, vol. 95, p. 045005, 2023.
\bibitem{temme2017} K. Temme, S. Bravyi, and J. M. Gambetta, ``Error mitigation for short-depth quantum circuits,'' \textit{Phys. Rev. Lett.}, vol. 119, p. 180509, 2017.
\bibitem{li2017} Y. Li and S. C. Benjamin, ``Efficient variational quantum simulator incorporating active error minimization,'' \textit{Phys. Rev. X}, vol. 7, p. 021050, 2017.
\bibitem{kandala2019} A. Kandala, K. Temme, A. D. C\'orcoles, A. Mezzacapo, J. M. Chow, and J. M. Gambetta, ``Error mitigation extends the computational reach of a noisy quantum processor,'' \textit{Nature}, vol. 567, pp. 491--495, 2019.
\bibitem{giurgica2020} T. Giurgica-Tiron, Y. Hindy, R. LaRose, A. Mari, and W. J. Zeng, ``Digital zero noise extrapolation for quantum error mitigation,'' in \textit{Proc. IEEE Int. Conf. Quantum Comput. Eng. (QCE)}, 2020, pp. 306--316.
\bibitem{wallman2016} J. J. Wallman and J. Emerson, ``Noise tailoring for scalable quantum computation via randomized compiling,'' \textit{Phys. Rev. A}, vol. 94, p. 052325, 2016.
\bibitem{kim2023} Y. Kim \textit{et al.}, ``Evidence for the utility of quantum computing before fault tolerance,'' \textit{Nature}, vol. 618, pp. 500--505, 2023.
\bibitem{klimov2018} P. V. Klimov \textit{et al.}, ``Fluctuations of energy-relaxation times in superconducting qubits,'' \textit{Phys. Rev. Lett.}, vol. 121, p. 090502, 2018.
\bibitem{proctor2020} T. Proctor, M. Revelle, E. Nielsen, K. Rudinger, D. Lobser, P. Maunz, R. Blume-Kohout, and K. Young, ``Detecting and tracking drift in quantum information processors,'' \textit{Nat. Commun.}, vol. 11, p. 5396, 2020.
\bibitem{rahman2026} R. Rahman and D. C. Nguyen, ``Quantum noise mitigation with adaptive zero-noise extrapolation: A contextual multi-armed bandits approach,'' \textit{IEEE J. Sel. Areas Commun.}, vol. 44, pp. 5600--5614, 2026, doi: 10.1109/JSAC.2026.3721807.
\bibitem{hour2026} L. Hour, M. Go, and Y. Han, ``Improving zero-noise extrapolation for quantum-gate error mitigation using a noise-aware folding method,'' \textit{IEEE Access}, vol. 14, pp. 37362--37372, 2026, doi: 10.1109/ACCESS.2026.3672056.
\bibitem{czarnik2021} P. Czarnik, A. Arrasmith, P. J. Coles, and L. Cincio, ``Error mitigation with Clifford quantum-circuit data,'' \textit{Quantum}, vol. 5, p. 592, 2021.
\bibitem{perezguijarro2026} J. P\'erez-Guijarro, A. Pag\`es-Zamora, and J. R. Fonollosa, ``Extension of Clifford data regression methods for quantum error mitigation,'' \textit{IEEE Trans. Quantum Eng.}, vol. 7, pp. 1--22, 2026, doi: 10.1109/TQE.2026.3677541.
\bibitem{prodius2026} M. Prodius, P. Czarnik, M. McKerns, A. T. Sornborger, and L. Cincio, ``Robust design under uncertainty in quantum error mitigation,'' \textit{IEEE Trans. Quantum Eng.}, vol. 7, pp. 1--13, 2026, doi: 10.1109/TQE.2026.3680641.
\bibitem{liao2024} H. Liao \textit{et al.}, ``Machine learning for practical quantum error mitigation,'' \textit{Nat. Mach. Intell.}, vol. 6, 2024.
\bibitem{placidi2026} L. Placidi, I. Williams, E. Rinaldi, D. Mills, C. C\^irstoiu, V. Eccles, and R. Duncan, ``Deep learning approaches to quantum error mitigation,'' arXiv:2601.14226, 2026.
\bibitem{li2010} L. Li, W. Chu, J. Langford, and R. E. Schapire, ``A contextual-bandit approach to personalized news article recommendation,'' in \textit{Proc. 19th Int. Conf. World Wide Web (WWW)}, 2010, pp. 661--670.
\bibitem{vaswani2017} A. Vaswani \textit{et al.}, ``Attention is all you need,'' in \textit{Advances in Neural Information Processing Systems}, vol. 30, 2017.
\bibitem{qiskit2024} A. Javadi-Abhari \textit{et al.}, ``Quantum computing with Qiskit,'' arXiv:2405.08810, 2024.
\bibitem{driftdata} ``qiskit-calibration-drift,'' Hugging Face dataset \texttt{phanerozoic/qiskit-calibration-drift}, revision 539bc99, 2026. [Online]. Available: \url{https://huggingface.co/datasets/phanerozoic/qiskit-calibration-drift}
\end{thebibliography}

%=====================================================================
\appendices
%=====================================================================
\section{Circuit Diagrams}
% Add your circuit images here (e.g. exported from Qiskit with qc.draw("mpl")).
\begin{figure}[htbp]
\centering
\figplaceholder{Circuit diagram: hea\_d4 (4 qubits, 12 CZ)}
\caption{Hardware-efficient ansatz, depth 4, after transpilation to \{CZ, sx, rz, x\}.}
\end{figure}
\begin{figure}[htbp]
\centering
\figplaceholder{Circuit diagram: tfim\_d3 and qaoa\_d2}
\caption{Trotterised TFIM (3 steps) and QAOA (2 layers) circuits.}
\end{figure}

\section{Experimental Settings}
\begin{table}[htbp]
\caption{Main settings}
\begin{center}
\begin{tabular}{p{2.3cm}p{5.6cm}}
\toprule
\textbf{Setting} & \textbf{Value} \\
\midrule
Software & Python 3.14.5, Qiskit 2.5.1, Aer 0.17.2, PyTorch 2.13.0, scikit-learn 1.9.1 \\
Window / horizons & 20 calibrations / 1, 2, 4, 7 \\
Transformer & 1 layer, width 32, 4 heads, FF 64, dropout 0.1, 5 seeds \\
Shot budget & 8,192 per observable per decision \\
ZNE scales & 1, 2, 3, 5, 7 (local CZ folding) \\
Twirling & 8 random Pauli frames per circuit \\
Coherent error (R2) & 0.03 rad controlled phase per CZ \\
Bootstrap & 10,000 cluster resamples, Holm correction \\
\bottomrule
\end{tabular}
\end{center}
\end{table}

\section{Code Listings}
The complete code is in the project repository; the following excerpts show the core steps.

\textit{Noise model from a calibration snapshot} (\texttt{src/qem/device.py}):
\begin{lstlisting}
for q in range(CHAIN_LEN):
    tr = _relax(state["t1"][q], state["t2"][q], sx_ns)
    p = _depol_fill(min(state["sx"][q], MAX_GATE_ERROR),
                    1 - average_gate_fidelity(tr), 2)
    err = _compose(tr, depolarizing_error(p, 1)) if p > 0 else tr
    nm.add_quantum_error(err, ["sx", "x"], [q])
\end{lstlisting}

\textit{Local CZ folding for ZNE} (\texttt{src/qem/mitigation.py}):
\begin{lstlisting}
def fold(ops, scale):
    n_cz = sum(o[0] == "cz" for o in ops)
    extra = int(round((scale - 1) * n_cz / 2))
    full, part = divmod(extra, n_cz)
    out, seen = [], 0
    for o in ops:
        out.append(o)
        if o[0] == "cz":
            out.extend([o, o] * (full + (1 if seen < part else 0)))
            seen += 1
    return out, (n_cz + 2 * extra) / n_cz
\end{lstlisting}

\textit{Model-based controller and LinUCB selection} (\texttt{src/qem/policies.py}):
\begin{lstlisting}
def controller_choice(pred):        # pred: predicted error per action
    return np.argmin(pred, axis=1)

def select(self, x):                # LinUCB, base paper Eqs. 17-18
    Ainv = np.linalg.inv(self.A)
    theta = np.einsum("kij,kj->ki", Ainv, self.b)
    bonus = np.sqrt(np.einsum("i,kij,j->k", x, Ainv, x))
    return int(np.argmax(theta @ x + self.eta * bonus
                         - self.lam_c * COST_NORM))
\end{lstlisting}

\section{Reproducibility}
All results can be regenerated with \texttt{python3 reproduce.py all}. The raw data are downloaded at a fixed revision and checked against a SHA-256 hash. Test evaluations refuse to run if any frozen input has changed (\texttt{configs/frozen\_*.yaml}). The pre-registered protocol, its change log and the stage reports are in the \texttt{docs} folder.

\end{document}
